\section{Alexandria and the Need for a Greek Torah}

Alexander's conquests did not make the eastern Mediterranean culturally uniform, but they made Greek the dominant language of administration and exchange across successor kingdoms. After Alexander's death in 323 \textsc{bc}, the Ptolemies ruled Egypt from Alexandria. The new capital joined royal government, a court-sponsored scholarly establishment, Mediterranean trade, Egyptian institutions, and communities drawn from many regions. Jews were established in Egypt before this period; under the Ptolemies they became a substantial Greek-speaking diaspora population whose life cannot be divided cleanly into “Jewish” and “Hellenistic” halves.

The Torah remained a Hebrew text and a marker of Israel's inherited law. Yet a community conducting public, commercial, educational, and domestic life in Greek needed forms of access that Hebrew alone could not provide for every member. Translation could teach, interpret, defend, and make the ancestral law intelligible. It could also acquire legal and civic significance for a community ordered by its own ancestral norms. None of those functions requires that Hebrew disappeared. The Greek Pentateuch's close engagement with Hebrew wording instead implies continuing reference to a source text whose verbal shape mattered.

\subsection{A royal library, a community project, or both?}

The \emph{Letter of Aristeas} makes the royal library the immediate cause: the librarian Demetrius advises Ptolemy II that the law of the Jews deserves inclusion, and the king secures a Jerusalem scroll and translators. Modern reconstructions offer other or additional motives---public reading, instruction, communal administration, apologetic prestige, or the needs of Greek-dominant Jews. The evidence does not permit one motive to eliminate all the others.

A royal-library setting is culturally intelligible in Ptolemaic Alexandria, where collection and textual scholarship served kingship. It is not therefore historically proved in the form narrated by Aristeas. Conversely, a communal need for Greek is a strong social inference, but no surviving meeting record identifies a synagogue committee or liturgical mandate. The most secure formulation is deliberately modest: a well-educated Jewish initiative in Ptolemaic Egypt, probably Alexandria, produced the Torah in Greek; the exact balance of patronage, law, worship, education, and cultural representation remains disputed.

\evidenceframe{Dirk Büchner's survey of the Greek Pentateuch and the linguistic character of the translation}{The impetus plausibly came from an educated Jewish sector still familiar with Hebrew; the Greek often keeps the source's words and structures in view while attempting varied solutions.}{Linguistic profile can illuminate translators and aims, but it cannot name a patron or reconstruct the institutional contract.}

\subsection{Scrolls before Bibles}

The first translation belonged to a scroll culture. Biblical books and groups of books circulated separately; the technological and social world of a single-volume Christian pandect lay centuries ahead. “The Hebrew Bible was translated” can therefore suggest an anachronistic source object. The Torah was a recognizable five-book collection, but the later Prophets, Writings, and other Jewish works did not arrive in Greek as one bound shipment.

Early Greek biblical artifacts confirm use while also enforcing humility. Rylands Greek Papyrus 458 preserves portions of Deuteronomy and is normally dated to the second century \textsc{bc}; it is a material witness to a Greek Torah not far removed from the originating period. Other papyri and the Greek Minor Prophets scroll from Na\d{h}al \d{H}ever show that Greek Jewish Scripture circulated and was revised before the great Christian codices. They do not preserve a complete Alexandrian edition or a fixed table of contents.

\begin{fourstable}{Evidence}{Approximate setting}{What it shows}{What it does not show}
Greek Torah language & Third-century \textsc{bc} translation layer & Jewish command of Hebrew and Greek, sustained engagement with source order and vocabulary. & The translators' names, exact workplace, or one exclusive motive.\\
Aristeas & Probably second-century \textsc{bc} Alexandrian Jewish narrative & A developed authorization story for the Greek Law and its public acceptance. & Eyewitness verification of Ptolemy's commission or the whole later corpus.\\
Rylands 458 & Second-century \textsc{bc} Deuteronomy fragments & Early physical circulation of Greek Torah. & A complete text, canon, or library copy.\\
Later codices & Fourth--fifth centuries \textsc{ad}, Christian & Large Greek Christian collections and important textual forms. & The exact contents or order of Jewish Scripture five or six centuries earlier.\\
\end{fourstable}

The translation belongs, then, inside Jewish cultural agency. Greek was not an alien solvent poured over an unchanged Hebrew object. Jewish translators made their law speak in the common language of their world, while the resistant texture of their Greek kept announcing that another language stood behind it.
